\scanpage{123}
On vérifie que $K'$ est encore quadratique totalement imaginaire sur un corps totalement réel $F'$, et que $\delta$ tue le sous-groupe algébrique $\operatorname{Ker}(F'{}^*\xrightarrow{N}\mathbb{G}_m)$ de $F'{}^*\subset K'{}^*$. Grâce à cela, si $\bar u\in\operatorname{Gal}(\overline{\mathbb{Q}}/K')^{ab}$ correspond à un idèle $\tilde u$ de $K'$ par le corps de classe, la composante finie de $\delta(\tilde u)$, $\in\prod_{v\nmid\infty}K_v^*$, ne dépend que de $u$, si prise modulo $K^*$.

\begin{tabular}{|p{0.92\linewidth}}
\textcircled{Th} Soit $(A,m,\alpha)$ du type dit. Pour $u\in\operatorname{Gal}(\overline{\mathbb{Q}}/K')$, la classe d'isomorphie de $u(A,m,\alpha)$ ne dépend que de l'image \deleted{de} $\bar u$ de $u$ dans $\operatorname{Gal}(\overline{\mathbb{Q}}/K')^{ab}$. C'est celle de $(A,m,\delta(\tilde u)^{-1}\circ\alpha)$
\end{tabular}

Exemple : $K$ quadratique imaginaire - On regarde des courbes elliptiques $A$ à multiplication complexe par l'anneau des entiers $O$ de $K$. On fixe un plongement $K\longrightarrow\overline{\mathbb{Q}}$, et on veut qu'il donne l'action sur Lie.

On a
\[
\begin{gathered}
\text{courbe ell. à mult. par }O\quad\longleftrightarrow\\[4pt]
\begin{gathered}
\text{courbe ell. à isogénie près, à mult. par }K,\text{ plus pour chaque }\ell\\
\text{un réseau }T_\ell\text{ dans }V_\ell(A),\text{ coïncidant pour presque tout }\ell\text{ avec le}\\
\text{réseau }\unsure{\text{standard}|\text{standart}}\text{ (qui n'a un sens que pour presque tout }\ell),\text{ ce réseau étant}\\
O\otimes\mathbb{Z}_\ell\text{-stable}
\end{gathered}\\
\longleftrightarrow\\
\begin{gathered}
\text{courbe ell à isogénie près, à multiplication par }K,\text{ plus pour chaque }\ell\\
\text{un isomorphisme }\alpha_\ell:V_\ell(A)\xrightarrow{\sim}K\otimes\mathbb{Q}_\ell,\text{ donné modulo }(O\otimes\mathbb{Z}_\ell)^*,\\
\text{et donnant un isomorphisme }V_{\unsure{\mathbb{A}^f|\mathbb{A}^+}}(A)\xrightarrow{\sim}\text{adèles finis de }K\\
\text{(un objet de type a) b) c), donné modulo }\widehat O^*\text{)}
\end{gathered}
\end{gathered}
\]
\margin{123}
